\chapter{Physical and Actuator Model}
\label{ch:physics}

The first prototype needed only gravity and thrust. That was enough to test the learning loop,
but every step toward a reusable-booster-shaped vehicle exposed another missing interaction.
Fuel changed the mass, grid fins required airflow, airflow required atmosphere and wind, and a
landing could not be judged without modelling individual feet and post-contact motion. The
physical model grew in that order. Its purpose is to create a coherent control problem, not to
claim flight-level fidelity.

\section{Rigid-body foundation and changing mass}

Unity integrates a six-degree-of-freedom rigid body. Let position be $\mathbf p$, velocity
$\mathbf v$, body angular velocity $\boldsymbol\omega$, mass $m$ and inertia tensor
$\mathbf I$. Translational and rotational motion conceptually follow
\begin{align}
  m\dot{\mathbf v} &= \mathbf F_T+\mathbf F_A+m\mathbf g, \\
  \mathbf I\dot{\boldsymbol\omega}+
  \boldsymbol\omega\times(\mathbf I\boldsymbol\omega) &=
  \boldsymbol\tau_T+\boldsymbol\tau_A+\boldsymbol\tau_R,
\end{align}
where the subscripts denote thrust, aerodynamic and RCS contributions. The simulator calculates
forces, application points, mass and inertia; Unity performs numerical integration and contact
resolution.

The Falcon-inspired reference body uses radius \SI{1.83}{m}, height \SI{41.2}{m}, dry mass
\SI{22200}{kg} and maximum propellant mass \SI{400000}{kg}. These are configurable nominal
values rather than a complete accounting of real vehicle hardware. The early cylinder kept a
constant mass, which meant the same throttle always produced the same acceleration. Once fuel
consumption was introduced, the combined centre of mass became
\begin{equation}
  \mathbf r_{\mathrm{cm}}=\frac{\sum_i m_i\mathbf r_i}{\sum_i m_i}.
\end{equation}
Approximate component inertias are combined with the parallel-axis theorem. Propellant behaves
as a changing rigid mass; liquid slosh is outside the model.

\section{From still air to a repeatable environment}

Grid fins made airflow unavoidable. The simulator first needed density and air-relative
velocity before it could calculate any useful aerodynamic control. Air density follows the
exponential approximation
\begin{equation}
  \rho(h)=\rho_0 k_\rho \exp\!\left(-\frac{h}{H}\right),
  \label{eq:density}
\end{equation}
where $h$ is altitude, $H$ is a scale height and $k_\rho$ is a user-selected multiplier. A
configurable horizontal wind vector is subtracted from vehicle velocity,
$\mathbf v_a=\mathbf v-\mathbf v_{\mathrm{wind}}$, and an optional low-order gust varies that
vector over time.

The intention is controlled disturbance, not weather simulation. Direction randomization and a
fixed seed make wind experiments repeatable, but the model has no altitude profile, turbulence
spectrum or correlated three-dimensional gust field. A wind setting therefore proves that the
force path exists; it does not by itself prove that a policy is wind-robust.

\section{Discovering that drag must work in both directions}

The first drag implementation looked plausible during descent but behaved incorrectly when the
rocket moved upward. The error came from losing the sign of axial velocity. The current model
decomposes air-relative velocity into axial and lateral components and writes axial drag as
\begin{equation}
  \mathbf F_{D,a}=-\tfrac12\rho C_{D,a}A_a
  v_a\lvert v_a\rvert\,\hat{\mathbf e}_a,
  \label{eq:axialdrag}
\end{equation}
with nominal $C_{D,a}=0.30$. The product $v_a\lvert v_a\rvert$ preserves direction, so the force
opposes both upward and downward motion.

Broadside drag uses nominal coefficient $C_{D,s}=1.0$, projected side area and lateral velocity.
It is applied at an approximate centre of pressure and therefore also produces torque. A small
base-drag term acts during tail-first descent. Separating these components made sign and scale
errors easier to test than a single global drag coefficient, although the result remains a
low-order model without Mach-dependent tables, compressibility or plume interaction.

\section{Grid fins and the limits of visual geometry}

Once body drag behaved consistently, the grid fins could contribute control forces. Their
response scales with dynamic pressure
\begin{equation}
  q=\tfrac12\rho\lVert\mathbf v_a\rVert^2.
\end{equation}
Each fin has a configurable envelope, deflection limit and slew rate. The nominal envelope is
\SI{1.59}{m} by \SI{1.23}{m}, with \SI{0.30}{m} thickness and an effective lift scale of 0.5.

An early version treated too much of that envelope as a solid aerodynamic plate. A real grid fin
contains open lattice, but resolving that flow is beyond the simulator. The model instead uses
the visible envelope for geometry and reduces its idealized response with the lift scale. This
does not reproduce lattice aerodynamics; it records the approximation explicitly and lets the
user vary it.

\section{Propulsion, fuel and actuator timing}

Adding more realistic engines changed the problem more than simply changing maximum thrust. An
engine command contains throttle and two gimbal coordinates. For each active engine,
\begin{equation}
  \mathbf F_T=T_{\max}\,k_T(h)\,u_T\,\hat{\mathbf e}_{\mathrm{gimbal}},
\end{equation}
where $u_T$ respects the configured minimum throttle and $k_T(h)$ approximates the change from
sea-level to vacuum thrust using the pressure ratio from Equation~\eqref{eq:density}. A sandbox
preset, for example, uses maximum sea-level thrust \SI{845}{kN}, specific impulse \SI{282}{s},
minimum throttle 0.39 and maximum gimbal angle \SI{7}{\degree}.

Fuel flow is calculated from
\begin{equation}
  \dot m_f=-\frac{T}{I_{\mathrm{sp}}g_0}.
  \label{eq:massflow}
\end{equation}
The effective $I_{\mathrm{sp}}$ is adjusted approximately with pressure. Neither this adjustment
nor the thrust interpolation is presented as a proprietary engine model.

The first engine responded instantly, allowing the policy to create control strategies that no
physical actuator could follow. Throttle spool, gimbal slew, startup delay, shutdown transient,
minimum run time and restart cooldown were introduced to close that gap. These constraints later
became part of the learning story themselves: strict timing made useful behaviour difficult to
discover in L8, while relaxed timing exposed excessive cycling and motivated L9--L11.

\section{From contact to a stable landing}

Touching the ground was initially enough to end an episode. That definition allowed one-foot
contacts, hard impacts and unstable rocking to appear successful. The final landing model exposes
four individual foot contacts and records first-contact velocity, tilt, angular rate, maximum
impulse, rebound height, stable-hold time, feet outside the pad and structural strikes. The
episode can therefore distinguish reaching the ground from completing a stable landing.

RCS channels follow the same modular pattern: they apply bounded forces and consume a separate
propellant budget using Equation~\eqref{eq:massflow}. The nominal thruster provides
\SI{750}{N}, has \SI{70}{s} specific impulse and respects a minimum \SI{0.10}{s} pulse.

Ground response still comes from Unity colliders and a general-purpose contact solver. The
simulator does not model crushing, detailed leg compliance, soil, structural deformation or
plume--ground interaction. A successful landing in this thesis therefore means success under
the documented collider geometry and thresholds, not proof that a real vehicle would survive.

\section{Verifying one layer before training the next}

Each addition created a new opportunity for motion that looked believable while being wrong. I
therefore used component-level tests for zero-wind drag direction, upward and downward motion,
vacuum and sea-level thrust scaling, fuel depletion, gimbal and throttle rates, RCS pulse timing,
centre-of-mass movement and leg-contact classification. These tests checked deterministic
properties before a policy was trained against them.

Training then acted as a second form of pressure testing. A policy repeats any profitable defect
far more consistently than a manual user. When a learned strategy violated the intended physics,
the behaviour was reduced to a component test and corrected there. This sequence---implement,
isolate, train and inspect---connected the physical model back to the development method in
Chapter~\ref{ch:scope}.
